% titel: Blatt 1, A2: Monoisotopische Elemente / A7: Isotope
% skript: Mortimer Kap. 2 (Isotope, Häufigkeit)

% Teil I: Werkzeuge, kein Lösungsweg.

Begriffe rund um Isotope: Was ein monoisotopisches Element ist, was
Isotopenhäufigkeit bedeutet und wann zwei Kerne \emph{keine} Isotope sind.

\begin{werkzeugkasten}[Isotope, Isotopenhäufigkeit, Reinelemente]
\begin{itemize}
  \item \textbf{Isotope} eines Elements haben dieselbe Protonenzahl $Z$, aber
        verschiedene Neutronenzahl $N$ (gleiche Ordnungszahl, verschiedene
        Massenzahl $A$).
  \item Zwei Kerne mit \emph{unterschiedlicher} Ordnungszahl sind
        \emph{keine} Isotope --- sie gehören zu verschiedenen Elementen.
  \item \textbf{Isotopenhäufigkeit} = relativer Anteil eines Isotops an allen
        Atomen des Elements (in Prozent). Die mittlere Atommasse/Molmasse ist
        das gewichtete Mittel über alle Isotope und ihre Häufigkeiten.
  \item Ein \textbf{Reinelement} (monoisotopisches Element) kommt in der
        Natur mit genau einem stabilen Isotop vor (z. B. Beryllium,
        Fluor, Natrium, Aluminium, Phosphor, Gold).
\end{itemize}
\textit{Skript:} Mortimer Kap. 2.4 (Isotope, Isotopenhäufigkeit), Übung 1.
\end{werkzeugkasten}

\begin{achtung}[Darauf achten]
\begin{itemize}
  \item Isotopie verlangt gleiche $Z$. Zwei verschiedene $Z$ $\Rightarrow$
        zwei verschiedene Elemente, keine Isotope.
  \item Mittlere (Mol-)Masse = $\sum_i (\text{Häufigkeit}_i \cdot
        \text{Masse}_i)$, Häufigkeiten als Bruch rechnen.
  \item Zahlen aus der Aufgabenstellung (Fe-Verteilung) übernehmen, nicht
        eigene Werte einsetzen.
\end{itemize}
\end{achtung}
